\documentclass{article}
\usepackage{hep-title}
\usepackage{titling}
\usepackage{siunitx}

% A horizontal rule above and below the series line, giving the
% classic preprint masthead a "lined" frame the plain titling package
% does not provide on its own.
\preseries{\begin{center}\rule{0.6\linewidth}{0.4pt}\\[4pt]\Large\scshape}
\postseries{\\[4pt]\rule{0.6\linewidth}{0.4pt}\end{center}}

\series{Whitmore Institute Preprint Series}
\arxiv{2611.04821}

\title{Anomalous Damping in Coupled Torsion Pendula}
\author[1]{E. Vasquez-Lindqvist\email{e.vasquez@example.edu}}
\author[1,2]{T. Okoye}
\affil[1]{Department of Applied Physics, Whitmore Institute}
\affil[2]{Now at the National Metrology Laboratory}
\date{\today}

\begin{document}

\maketitle

\begin{abstract}
We report an anomalous damping term in a pair of magnetically
coupled torsion pendula that cannot be explained by eddy-current
losses alone. The excess damping scales with the square of coupling
strength rather than linearly, suggesting a parametric loss channel
not captured by the standard coupled-oscillator model.
\end{abstract}

\section{Introduction}

Coupled torsion pendula are a standard testbed for precision
measurements of weak forces, and their damping behaviour is normally
well described by independent eddy-current losses in each pendulum
plus a coupling-independent background. In the course of a routine
calibration run, we observed a damping excess that grows with
coupling strength in a way the standard model does not predict.

\section{Apparatus}

Two aluminium torsion pendula, each with a natural period of
\SI{4.2}{\second}, were coupled via a pair of neodymium magnets
mounted at matching radii on each pendulum's disk. Coupling strength
was varied by adjusting the gap between the magnet pairs across
seven settings, from \SI{2}{\milli\meter} (strongest coupling) to
\SI{14}{\milli\meter} (weakest).

\section{Results}

At each coupling setting, the free-decay envelope of both pendula was
fitted to an exponential to extract a quality factor $Q$. The
standard model predicts $Q^{-1}$ should be independent of coupling
strength, since eddy-current losses in each pendulum's aluminium disk
do not depend on the neighbouring pendulum. Instead, we observe
$Q^{-1}$ growing approximately quadratically with the inverse of the
magnet gap, i.e. with the square of coupling strength.

We considered three conventional explanations before concluding the
loss channel is genuinely anomalous: (i) mutual eddy-current coupling
between the two disks, which would predict a linear rather than
quadratic scaling; (ii) mechanical crosstalk through the shared
mounting frame, ruled out by repeating the measurement with the
pendula on independent mounts; and (iii) air damping changes from the
closer magnet proximity, ruled out by repeating in vacuum with no
change to the scaling exponent.

\section{Discussion}

The quadratic scaling is consistent with a parametric loss mechanism
in which the coupling itself modulates an internal degree of freedom
in each magnet's mounting, rather than acting purely as a
lossless torque coupling between the two pendula. We do not yet have
a first-principles model for this channel, and flag it as an open
question for follow-up work with magnet mounts of varying stiffness.

\end{document}
